% PROVISIONAL DRAFT -- not for submission. Every empirical number is a \Num... macro generated from
% committed reports (paper/generated/numbers.tex). Sections marked [PENDING] depend on data or
% authorization that does not exist yet; they must not be filled with unmeasured values.
\documentclass[sigconf,review,anonymous]{acmart}
% Generated by `uv run python -m www_paper.macros`; do not edit by hand.
\newcommand{\NumUsdcMedAbsPre}{2.06} % reports/c3/fragmentation_persistence.json: wedges/usdc/all_legs_traded/pre_crisis/median_abs_bps
\newcommand{\NumUsdcMedAbsCrisis}{8.13} % reports/c3/fragmentation_persistence.json: wedges/usdc/all_legs_traded/crisis/median_abs_bps
\newcommand{\NumUsdcMedAbsPost}{5.10} % reports/c3/fragmentation_persistence.json: wedges/usdc/all_legs_traded/post/median_abs_bps
\newcommand{\NumUsdcPNinetyFivePre}{8.6} % reports/c3/fragmentation_persistence.json: wedges/usdc/all_legs_traded/pre_crisis/p95_abs_bps
\newcommand{\NumUsdcPNinetyFiveCrisis}{77.8} % reports/c3/fragmentation_persistence.json: wedges/usdc/all_legs_traded/crisis/p95_abs_bps
\newcommand{\NumUsdcShareAboveTwentyFiveCrisis}{21.9} % reports/c3/fragmentation_persistence.json: wedges/usdc/all_legs_traded/crisis/thresholds/25.0/share_observed_minutes_above
\newcommand{\NumUsdcShareAboveTwentyFivePre}{0.3} % reports/c3/fragmentation_persistence.json: wedges/usdc/all_legs_traded/pre_crisis/thresholds/25.0/share_observed_minutes_above
\newcommand{\NumUsdcCrisisMinutes}{2,570} % reports/c3/fragmentation_persistence.json: wedges/usdc/all_legs_traded/crisis/observed_minutes
\newcommand{\NumUsdcHalfLifeOneMinCrisis}{0.5} % reports/c3/fragmentation_persistence.json: wedges/usdc/all_legs_traded/crisis/half_life_to_parity_minutes
\newcommand{\NumUsdtPNinetyFiveCrisis}{7.7} % reports/c3/fragmentation_persistence.json: wedges/usdt/all_legs_traded/crisis/p95_abs_bps
\newcommand{\NumUsdcDiffMedAbs}{6.08} % reports/c3/robustness.json: regime_differences/usdc/crisis_minus_pre_crisis/median_abs/estimate
\newcommand{\NumUsdcDiffMedAbsLo}{3.75} % reports/c3/robustness.json: regime_differences/usdc/crisis_minus_pre_crisis/median_abs/ci95/0
\newcommand{\NumUsdcDiffMedAbsHi}{9.16} % reports/c3/robustness.json: regime_differences/usdc/crisis_minus_pre_crisis/median_abs/ci95/1
\newcommand{\NumUsdcHalfLifeFifteenCrisis}{12.6} % reports/c3/robustness.json: horizons/usdc/15/crisis/half_life_to_parity_minutes
\newcommand{\NumUsdcHalfLifeSixtyCrisis}{16.2} % reports/c3/robustness.json: horizons/usdc/60/crisis/half_life_to_parity_minutes
\newcommand{\NumUsdcCoinbaseMedAbsCrisis}{10.78} % reports/c3/robustness.json: alternative_dollar_leg/usdc_coinbase_btcusd/crisis/median_abs_bps
\newcommand{\NumAnchorExtremeBps}{-1,470} % reports/events/E2023-03-USDC-SVB_protocol.json: r1/extreme_deviation_bps
\newcommand{\NumAnchorBreachMinutes}{1,379} % reports/events/E2023-03-USDC-SVB_protocol.json: r1/max_breach_minutes_in_window
\newcommand{\NumAnchorEventMedAbs}{209.2} % reports/events/E2023-03-USDC-SVB_protocol.json: windows/event/diagnostics/median_abs_dev_bps
\newcommand{\NumAnchorEventPNinetyFive}{950.8} % reports/events/E2023-03-USDC-SVB_protocol.json: windows/event/diagnostics/p95_abs_dev_bps
\newcommand{\NumAnchorEventShareHundred}{60} % reports/events/E2023-03-USDC-SVB_protocol.json: windows/event/diagnostics/share_traded_minutes_ge_100bps
\newcommand{\NumAnchorEventHalfLife}{1,212} % reports/events/E2023-03-USDC-SVB_protocol.json: windows/event/diagnostics/half_life_to_parity_minutes_15m_bins
\newcommand{\NumAnchorBaselineMedAbs}{1.0} % reports/events/E2023-03-USDC-SVB_protocol.json: windows/baseline/diagnostics/median_abs_dev_bps
\newcommand{\NumAnchorRecoveryMedAbs}{5.0} % reports/events/E2023-03-USDC-SVB_protocol.json: windows/recovery/diagnostics/median_abs_dev_bps
\newcommand{\NumSampleRows}{7,264,627} % reports/c1/march01_sample_validation.json: rows
\newcommand{\NumAnchorDvOne}{208} % reports/events/E2023-03-USDC-SVB_protocol.json: plan/dv1/estimate
\newcommand{\NumAnchorDvOneLo}{92} % reports/events/E2023-03-USDC-SVB_protocol.json: plan/dv1/ci95/0
\newcommand{\NumAnchorDvOneHi}{336} % reports/events/E2023-03-USDC-SVB_protocol.json: plan/dv1/ci95/1
\newcommand{\NumAnchorRecoveryHours}{68} % reports/events/E2023-03-USDC-SVB_protocol.json: plan/dv3/recovery_minutes
\newcommand{\NumAnchorBtcVolBaseline}{40} % reports/events/E2023-03-USDC-SVB_protocol.json: plan/covariates/baseline/btc_realized_vol_bps_per_hour
\newcommand{\NumAnchorBtcVolEvent}{80} % reports/events/E2023-03-USDC-SVB_protocol.json: plan/covariates/event/btc_realized_vol_bps_per_hour
\newcommand{\NumAnchorDvFour}{6.5} % reports/events/E2023-03-USDC-SVB_protocol.json: plan/dv4/estimate
\newcommand{\NumAnchorDvFourLo}{4.6} % reports/events/E2023-03-USDC-SVB_protocol.json: plan/dv4/ci95/0
\newcommand{\NumAnchorDvFourHi}{9.4} % reports/events/E2023-03-USDC-SVB_protocol.json: plan/dv4/ci95/1
\newcommand{\NumAnchorDvFive}{95} % reports/events/E2023-03-USDC-SVB_protocol.json: plan/dv5/ratio
\newcommand{\NumFdusdExtreme}{-1,279} % reports/events/E2025-04-FDUSD_protocol.json: r1/extreme_deviation_bps
\newcommand{\NumFdusdDvOne}{45} % reports/events/E2025-04-FDUSD_protocol.json: plan/dv1/estimate
\newcommand{\NumFdusdDvOneLo}{34} % reports/events/E2025-04-FDUSD_protocol.json: plan/dv1/ci95/0
\newcommand{\NumFdusdDvOneHi}{51} % reports/events/E2025-04-FDUSD_protocol.json: plan/dv1/ci95/1
\newcommand{\NumFdusdDvFour}{0.02} % reports/events/E2025-04-FDUSD_protocol.json: plan/dv4/estimate
\newcommand{\NumFdusdDvFourLo}{-0.02} % reports/events/E2025-04-FDUSD_protocol.json: plan/dv4/ci95/0
\newcommand{\NumFdusdDvFourHi}{0.05} % reports/events/E2025-04-FDUSD_protocol.json: plan/dv4/ci95/1
\newcommand{\NumFdusdRecoveryHours}{130} % reports/events/E2025-04-FDUSD_protocol.json: plan/dv3/recovery_minutes
\newcommand{\NumFdusdVolControl}{91} % reports/events/E2025-04-FDUSD_protocol.json: plan/covariates/control/btc_realized_vol_bps_per_hour
\newcommand{\NumFdusdVolEvent}{69} % reports/events/E2025-04-FDUSD_protocol.json: plan/covariates/event/btc_realized_vol_bps_per_hour
\newcommand{\NumPlaceboDvOne}{-3} % reports/events/P2023-02-BUSD-PAXOS-HALT_protocol.json: plan/dv1/estimate
\newcommand{\NumPlaceboDvFive}{5.4} % reports/events/P2023-02-BUSD-PAXOS-HALT_protocol.json: plan/dv5/ratio
\newcommand{\NumNegativeExtreme}{361} % reports/events/N2023-06-USDT-3POOL_robustness_binanceus.json: r1/extreme_deviation_bps
\newcommand{\NumNegativeExtremeBitstamp}{-40} % reports/events/N2023-06-USDT-3POOL_protocol.json: r1/extreme_deviation_bps
\newcommand{\NumTusdExtreme}{-297} % reports/events/E2024-01-TUSD_protocol.json: r1/extreme_deviation_bps
\newcommand{\NumUstExtreme}{-32,190} % reports/events/E2022-05-UST-COLLAPSE_protocol.json: r1/extreme_deviation_bps
\newcommand{\NumTusdDvOne}{113} % reports/events/E2024-01-TUSD_protocol.json: plan/dv1/estimate
\newcommand{\NumTusdDvOneLo}{100} % reports/events/E2024-01-TUSD_protocol.json: plan/dv1/ci95/0
\newcommand{\NumTusdDvOneHi}{121} % reports/events/E2024-01-TUSD_protocol.json: plan/dv1/ci95/1
\newcommand{\NumTusdDvFour}{-0.08} % reports/events/E2024-01-TUSD_protocol.json: plan/dv4/estimate
\newcommand{\NumUsdeExtreme}{-4,080} % reports/events/E2025-10-USDE-BINANCE_protocol.json: r1/extreme_deviation_bps
\newcommand{\NumUsdeUsdtGate}{15} % reports/events/E2025-10-USDE-BINANCE_protocol.json: plan/usdt_gate/median_abs_bps
\newcommand{\NumRegistryEntries}{16} % configs/event_registry.json
\newcommand{\NumRegistryClusters}{13} % configs/event_registry.json
\newcommand{\NumAcceptedClusters}{4} % configs/event_registry.json
\newcommand{\NumDocumentedPrimaries}{2} % configs/event_registry.json
\newcommand{\NumUnresolved}{3} % configs/event_registry.json


\begin{document}

\title{When Arbitrage Cannot Restore the Peg: Executable Capacity and Stablecoin Fragmentation Across Stress Events}

\begin{abstract}
Stablecoins are meant to trade at par across venues and quote currencies, yet stress episodes repeatedly
produce prices that differ across markets for hours or days. We ask whether limits on \emph{executable}
arbitrage capacity---depth, redemption access, settlement timing and inventory---help explain how large such
fragmentation becomes, how long it lasts, and how it resolves. We build a pre-registered multi-event design:
a sourced registry of \NumRegistryEntries{} candidate episodes from 2020--2026 in \NumRegistryClusters{}
independence clusters, explicit severity and independence rules, matched control windows, placebo and
expected-negative cases, and events after 2024 held out for out-of-sample evaluation. [PENDING: cross-event
results.] In the anchor episode (USDC, March 2023), the stablecoin's own dollar price stayed dislocated with a
half-life of about \NumAnchorEventHalfLife{} minutes, while within-venue cross-quote wedges reverted within
minutes, which separates redemption-market dislocation from cross-quote arbitrage failure. All results are
descriptive; executable capacity is not yet measured.
\end{abstract}

\maketitle

\section{Introduction}
Fiat-backed and synthetic stablecoins rely on arbitrage to hold the peg: when the secondary-market price
falls below par, traders buy and redeem; when relative prices across quote currencies diverge, traders route
across them. Both mechanisms need \emph{capacity}---the quantity that can actually be executed at a given cost
within a given time. Prior work shows that arbitrage design matters for peg stability~\cite{lyons2023keeps},
that primary-market access shapes crisis dynamics beyond what secondary prices reveal~\cite{watsky2024primary,du2025shadow},
and that cross-exchange crypto price deviations can be large and persistent when arbitrage capital is
constrained~\cite{makarov2020trading}.

\paragraph{Why multiple events.} A single crisis cannot separate mechanisms: in March 2023, bank-reserve risk,
weekend redemption closure and exchange-specific liquidity all moved together. We therefore pre-specify a
multi-event design covering algorithmic collapses, bank-reserve runs, exchange-failure contagion, issuer
solvency allegations and venue-infrastructure failures, with negative cases and placebo shocks evaluated
under the same rules. This improves on a March 2023-only study in three ways: (i) heterogeneity across
mechanisms becomes testable rather than assumed away; (ii) selection is rule-based and includes events
expected \emph{not} to qualify; (iii) post-2024 events are held out, so conclusions drawn from the in-sample
period face an honest test.

\paragraph{Contributions (current status).}
(1) A sourced, machine-validated event registry with independence clustering and coverage matrices
(\S\ref{sec:events}). (2) A measurement framework that separates peg dislocation, within-venue triangle wedges and
cross-venue wedges, and separates parity deviations from gross, executable and realized arbitrage
(\S\ref{sec:framework}). (3) Anchor-event evidence that these objects behave very differently
(\S\ref{sec:anchor}). (4) [PENDING] Cross-event comparisons and executable-capacity measurement.

\section{Related Work}
\label{sec:related}
Arbitrage and segmentation in crypto markets are documented by~\cite{makarov2020trading}. Stablecoin peg
mechanics and the role of arbitrage access are studied by~\cite{lyons2023keeps}; the legal and run-risk
framing by~\cite{zhang2023taming}; primary versus secondary markets and the SVB episode by
Federal Reserve staff~\cite{watsky2024primary,du2025shadow}. [PENDING: a systematic literature review, restricted
to verified sources, covering limits to arbitrage, market microstructure of crypto order books, AMM
liquidity, and stablecoin event studies.]

\section{Measurement Framework}
\label{sec:framework}
For stablecoin $q$ with USD peg, let $S_t^{q}$ be its USD price on a venue. We distinguish:
\begin{itemize}
\item \textbf{Peg dislocation} $D_t^{q} = 10^4 \log S_t^{q}$ (bps).
\item \textbf{Within-venue triangle wedge} $W_t^{q} = 10^4[\log P_t(\mathrm{BTC}/q) + \log S_t^{q} - \log P_t(\mathrm{BTC}/\mathrm{USD})]$,
which is zero under frictionless cross-quote arbitrage on one venue regardless of $S_t^q$.
\item \textbf{Cross-venue wedge}: the same construction with legs from different venues.
\end{itemize}
A nonzero $D$ or $W$ is a \emph{parity deviation}. It is a \emph{gross opportunity} only after quoted
spreads, a \emph{executable opportunity} only after fees, depth, latency, settlement and inventory, and a
\emph{realized trade} only when observed in trade data. Restoration capacity $K_t(\delta)$ is the
largest quantity that can be executed around the loop at all-in cost at most $\delta$ within horizon $h$.
Candle data identify parity deviations only; $K_t$ requires order-book or on-chain execution state.

\section{Hypotheses (pre-specified)}
\textbf{H1.} Fragmentation magnitude and persistence increase in stress windows relative to matched controls.
\textbf{H2.} Peg dislocations are more persistent than within-venue wedges when redemption access is impaired.
\textbf{H3.} Conditional on stress, lower executable restoration capacity is associated with larger and more persistent wedges.
H1--H2 are testable with candle data; H3 requires validated execution state and is [PENDING].

\section{Event Selection}
\label{sec:events}
Rules were fixed before measuring any non-anchor event. \emph{R1 (severity):} at least 30 traded minutes with
$|D| \ge 100$ bps inside some 24-hour window; $\ge 500$ bps marks the severe class. \emph{R2 (independence):} one
episode per primary trigger; contagion legs share their primary's cluster and never add episodes.
\emph{R3 (windows):} 7-day baseline ending two days before the event, the event window, a 7-day recovery, and a
7-day matched control ending 28 days before. \emph{R4 (sourcing):} at least two distinct publishers.
Events from 2024 onward are out-of-sample.

The registry holds \NumRegistryEntries{} entries in \NumRegistryClusters{} clusters:
\NumDocumentedPrimaries{} source-documented primary candidates (including one placebo---the BUSD issuance
halt~\cite{paxos2023busd}---and one expected negative case), \NumUnresolved{} unresolved candidates, and
\NumAcceptedClusters{} accepted cluster so far. [PENDING: R1 measurement for UST~\cite{jump2022ust}, TUSD, FDUSD,
USDe and the placebo; the target of at least three independent accepted episodes is \emph{not yet met}.]

\section{Data}
Anchor data: a 30,240-minute Binance.US one-minute candle panel (2023-03-01 to 2023-03-21), hash-pinned and
audited. Order-book reconstruction is validated on one public normalized L2 sample
(\NumSampleRows{} rows; BTCUSD, 2023-03-01)~\cite{tardis2026binanceus}; normalized feeds omit native exchange
sequence identifiers~\cite{coinapi2026limitbook}, so continuity is disclosed as unverifiable. Multi-event
candle data are listed in public archives but not yet acquired. Candle closes are last trades, not quotes.

\section{Anchor Event: USDC, March 2023 (provisional)}
\label{sec:anchor}
\textbf{Peg dislocation.} On Binance.US, R1 is met with \NumAnchorBreachMinutes{} breaching minutes in the worst
24-hour window and an extreme of \NumAnchorExtremeBps{} bps. The median $|D|$ was \NumAnchorBaselineMedAbs{} bps in
the baseline, \NumAnchorEventMedAbs{} bps in the event window (\NumAnchorEventShareHundred\% of traded minutes
beyond 100 bps) and \NumAnchorRecoveryMedAbs{} bps in recovery. On 15-minute medians the event-window half-life
to par was about \NumAnchorEventHalfLife{} minutes.

\textbf{Triangle wedge.} The median $|W^{\mathrm{USDC}}|$ rose from \NumUsdcMedAbsPre{} to \NumUsdcMedAbsCrisis{} bps
(difference \NumUsdcDiffMedAbs{} bps, 95\% block-bootstrap CI [\NumUsdcDiffMedAbsLo, \NumUsdcDiffMedAbsHi]); its 95th
percentile rose from \NumUsdcPNinetyFivePre{} to \NumUsdcPNinetyFiveCrisis{} bps, and \NumUsdcShareAboveTwentyFiveCrisis\% of
jointly traded crisis minutes exceeded 25 bps (\NumUsdcShareAboveTwentyFivePre\% before). USDT wedges stayed small
(crisis P95 \NumUsdtPNinetyFiveCrisis{} bps). Yet $W$ reverted fast: half-life \NumUsdcHalfLifeOneMinCrisis{} minutes on
one-minute closes and \NumUsdcHalfLifeFifteenCrisis--\NumUsdcHalfLifeSixtyCrisis{} minutes on 15--60-minute medians. Using
Coinbase BTC-USD as the dollar leg gives a crisis median of \NumUsdcCoinbaseMedAbsCrisis{} bps.

\textbf{Interpretation.} Cross-quote arbitrage kept relative prices close while USDC's dollar value stayed
dislocated---consistent with H2 and with impaired redemption~\cite{du2025shadow}, but one event cannot
establish it.

\section{Executable Capacity [PENDING]}
Capacity estimation is gated on validated crisis-period L2 or on-chain execution state and explicit
authorization. No capacity figures are reported.

\section{Limitations}
Single measured event; last-trade candles with asynchronous legs and bid-ask bounce; calendar regime labels;
USDT-quoted archives require USD correction; no executable depth; normalized L2 lacks native sequence IDs.

\section{Reproducibility}
Code, registry, rules and reports are versioned; raw licensed data are referenced by hash. Every number in
this draft is generated by \texttt{python -m www\_paper.macros} from committed JSON reports.

\bibliographystyle{ACM-Reference-Format}
\bibliography{references}
\end{document}
